% ===== §4 =====
\section{Symbol and Interpretation}
\label{sec:hermeneutic}

\noindent
The second tradition is the one this paper most nearly resembles and must therefore most carefully distinguish itself from, for it comes closest to the truth and misses it in the most instructive way. Its thesis is that culture is meaning, and that the medium of meaning is the symbol. Cassirer gave the thesis its widest philosophical frame: the human being is not so much a rational animal as a symbolic one, an \emph{animal symbolicum}, and culture is the totality of the symbolic forms---language, myth, art, religion, science---through which a world is constituted at all rather than merely encountered \citep{cassirer1944,cassirer1953}. On this view culture is not one region of human life but the medium in which human life as such takes place, since there is no access to a world except through some symbolic form that shapes it. Geertz brought the thesis down from the philosophy of forms to the practice of interpretation. Believing, with Weber, that the human being is an animal suspended in webs of significance he himself has spun, Geertz took culture to be those webs, and the study of it to be not an experimental science in search of law but an interpretive one in search of meaning \citep{geertz1973}. Culture is public because meaning is public: it lives in symbols that can be read, and the reading, his \emph{thick description}, is the patient recovery of the layered significance that separates a wink from a twitch, a significance that the physical motion alone cannot carry and that only the web of meaning supplies.

This is a great advance over the first tradition, and along the axis this paper cares about most it moves in the right direction. It sees that culture is not a stock of items but a texture of meaning; that meaning is relational, a matter of a sign's place in a web rather than of any property it carries in itself; and that culture is therefore closer to a language than to a patrimony. The metaphor of the web is almost the metaphor this paper wants, for a web is a structure of relations and not a heap of things. And Geertz's insistence that culture is spun by those suspended in it seems, for a moment, to grant exactly the generativity the first tradition denied: the web is woven, and weaving is an act.

\subsection*{Where the web hangs already spun}

But the moment does not hold, and the reason it does not hold is the presupposition this tradition shares, beneath its advances, with the one before it. The web, in the end, is encountered as \emph{already spun}. For all that Geertz says the webs are spun by those suspended in them, the suspension is prior in his account to any spinning one could observe: the individual is born into webs of significance that precede him, and interpretation, thick description, is addressed to meaning that is \emph{already there} to be read, a public text the ethnographer arrives late to and construes. The generativity in the metaphor is grammatically present and analytically absent. One never, in this tradition, catches the web in the act of being woven; one finds it woven and reads it. And Cassirer's symbolic forms are still more nearly a priori: they are the standing conditions of any experience whatever, the forms a world must take to be a world, and as such they precede every particular act of meaning as its condition rather than issuing from it. In the terms of the map, the tradition has moved culture toward the relational and the structural in its account of the \emph{bearer} of meaning---meaning lives in the web, between the signs---while leaving its \emph{mode of being} settled: the web is a prior structure into which the individual is inducted, and meaning is found, decoded, interpreted, never watched as it comes to be. The symbolic system is treated as anterior to the subjects who move within it, and the subject is figured as one who enters a language already spoken.

\subsection*{The symbolic order and the lack it opens}

It is exactly here, at the point where the symbolic system is granted priority over the subject, that the psychoanalytic line of this paper first enters, and it enters not to refute this tradition but to name a dimension the tradition leaves in the dark. For if the subject is inducted into a symbolic order that precedes him, then the relation between the subject and that order is not neutral, and psychoanalysis is the discourse that has taken the measure of what that induction costs. In the Lacanian account the symbolic order is the register of language and law, the treasury of signifiers that is there before the subject and that the subject must enter to become a speaking being at all; and this entry is not a gain only. To enter language is to submit what one is to what can be said of it, and since no signifier can ever say the whole of what the subject is, the entry opens, at the very heart of the subject, a constitutive lack, an emptiness that the symbolic order institutes by the same act with which it constitutes the subject \citep{lacan2006ecrits}. The symbolic order that Cassirer and Geertz describe from the outside, as the medium of meaning, is described by psychoanalysis from the inside, as the big Other into which desire is inducted and by which it is marked with a lack it can never fill. This is the first appearance of a thesis the paper will develop at length: that culture, considered as a symbolic system, is not a neutral web of meaning but the very apparatus by which drive is first coded, and that the coding leaves a remainder. The interpretive tradition sees the web and reads its meanings; it does not see that the web is stretched over a hole, that the subject suspended in significance is suspended over its own constitutive lack, and that the meanings it spins are, among other things, ways of living beside that lack. What the tradition cannot see, because it stays at the level of the symbolic taken alone, is precisely what the later sections of this paper must recover: that culture spans more than the symbolic, and that beneath the web of meaning runs a register the web does not exhaust.
